\subsection{DEFINITION OF NILPOTENT LIE ALGEBRAS}

DEFINITION 1. A Lie algebra \(\mathfrak{g}\) is called nilpotent if there exists a decreasing finite sequence of ideals \((\mathfrak{g}_i)_{1\leqslant i\leqslant p}\) of \(\mathfrak{g}\) with \(\mathfrak{g}_0 = \mathfrak{g},\mathfrak{g}_p = \{0\}\) , such that \([\mathfrak{g},\mathfrak{g}_i]\subset \mathfrak{g}_{i + 1}\) for \(0\leqslant i < p\) .

A commutative Lie algebra is nilpotent.

PROPOSITION 1. Let \(\mathfrak{g}\) be a Lie algebra. The following conditions are equivalent:

\begin{enumerate}
\def\labelenumi{(\alph{enumi})}
\item
  g is nilpotent;
\item
  \(\mathcal{C}^k\mathfrak{g} = \{0\}\) for sufficiently large \(k\) ;
\item
  \(\mathcal{C}_k\mathfrak{g} = \mathfrak{g}\) for sufficiently large \(k\) ;
\item
  there exists an integer \(k\) such that \(\operatorname{ad} x_1 \circ \operatorname{ad} x_2 \circ \dots \circ \operatorname{ad} x_k = 0\) for all elements \(x_1, x_2, \ldots, x_k\) in \(\mathfrak{g}\) ;
\item
  there exists a decreasing sequence of ideals \((\mathfrak{g}_i)_{0\leqslant i\leqslant n}\) of \(\mathfrak{g}\) with \(\mathfrak{g}_0 = \mathfrak{g},\mathfrak{g}_n = \{0\}\) , such that \([\mathfrak{g},\mathfrak{g}_i]\subset \mathfrak{g}_{i + 1}\) and \(\dim \mathfrak{g}_i / \mathfrak{g}_{i + 1} = 1\) for \(0\leqslant i <   n\)
\end{enumerate}

If \(\mathcal{C}^k\mathfrak{g} = \{0\}\) (resp. \(\mathcal{C}_k\mathfrak{g} = \mathfrak{g}\) ), clearly the sequence \(\mathcal{C}^1\mathfrak{g},\ldots ,\mathcal{C}^k\mathfrak{g}\) (resp. \(\mathcal{C}_k\mathfrak{g},\mathcal{C}_{k - 1}\mathfrak{g},\ldots ,\mathcal{C}_0\mathfrak{g}\) ) has the properties of Definition 1 and hence \(\mathfrak{g}\) is nilpotent. Conversely, suppose that there exists a sequence \((\mathfrak{g}_i)_{0\leqslant i\leqslant p}\) with the properties of Definition 1. It is seen by induction on \(i\) that \(\mathfrak{g}_i\supset \mathcal{C}^{i + 1}\mathfrak{g}\) and \(\mathfrak{g}_{p - i}\subset \mathcal{C}_i\mathfrak{g}\) . Hence \(\mathcal{C}^{p + 1}\mathfrak{g} = \{0\}\) and \(\mathcal{C}_p\mathfrak{g} = \mathfrak{g}\) . We have thus proved that conditions (a), (b) and (c) are equivalent. On the other hand, \(\mathcal{C}^i\mathfrak{g}\) is the set of linear combinations of elements of the form

\[
\left[ x _ {1}, \left[ x _ {2}, \dots , \left[ x _ {i - 2}, \left[ x _ {i - 1}, x _ {i} \right] \right] \dots \right] \right]
\]

where \(x_{1}, x_{2}, \ldots, x_{i}\) run through \(\mathfrak{g}\) . Hence conditions (b) and (d) are equivalent. Finally, if there exists a sequence \((\mathfrak{g}_i)_{0 \leqslant i \leqslant p}\) of ideals with the properties of

Definition 1, there exists a decreasing sequence \((\mathfrak{h}_i)_{0\leqslant i\leqslant n}\) of vector subspaces of \(\mathfrak{g}\) of dimensions \(n,n - 1,n - 2,\ldots ,0\) and a sequence of indices

\[
i _ {0} <   i _ {1} <   \dots <   i _ {p}
\]

with \(\mathfrak{g}_0 = \mathfrak{h}_{i_0},\mathfrak{g}_1 = \mathfrak{h}_{i_1}\dots ,\mathfrak{g}_p = \mathfrak{h}_{i_p}\) ; then as \([\mathfrak{g},\mathfrak{h}_{i_k}]\subset \mathfrak{h}_{i_{k + 1}}\) the \(\mathfrak{h}_i\) are ideals and \([\mathfrak{g},\mathfrak{h}_i]\subset \mathfrak{h}_{i + 1}\) for all \(i\) . Hence conditions (a) and (e) are equivalent.

COROLLARY 1. The centre of a non-zero nilpotent Lie algebra is non-zero.

COROLLARY 2. The Killing form of a nilpotent Lie algebra is zero.

For all \(x\) and \(y\) in a nilpotent Lie algebra ad \(x \circ\) ad \(y\) is nilpotent and hence of zero trace.

PROPOSITION 2. Subalgebras, quotient algebras and central extensions of a nilpotent Lie algebra are nilpotent. A finite product of nilpotent Lie algebras is a nilpotent Lie algebra.

Let g be a Lie algebra, g' a subalgebra of g, b an ideal of g, t = g/b and φ the canonical mapping of g onto t. If g is nilpotent, then \(C^{k}g = \{0\}\) for some integer k, hence \(C^{k}g' \subset C^{k}g = \{0\}\) and \(C^{k}t = \phi(C^{k}g) = \{0\}\) and hence g' and t are nilpotent. If t is nilpotent and b is contained in the centre of g, then \(C^{k}t = \{0\}\) for some integer k, hence \(C^{k}g \subset b\) and therefore \(C^{k+1}g \subset [b, g] = \{0\}\) , so that g is nilpotent. Finally, the assertion concerning products follows for example from the assertion (a) \(\Leftrightarrow\) (d) of Proposition 1.

Definition 1 and Proposition 2 show that nilpotent Lie algebras are precisely the algebras obtained from commutative Lie algebras by a sequence of central extensions.

PROPOSITION 3. Let g be a nilpotent Lie algebra and h a subalgebra of g distinct from g. The normalizer of h in g is distinct from h.

Let \(k\) be the greatest integer such that \(\mathcal{C}^k\mathfrak{g} + \mathfrak{h}\neq \mathfrak{h}\) . Then

\[
[ \mathcal {C} ^ {k} \mathfrak {g} + \mathfrak {h}, \mathfrak {h} ] \subset \mathcal {C} ^ {k + 1} \mathfrak {g} + \mathfrak {h} \subset \mathfrak {h}
\]

and hence the normalizer of \(\mathfrak{h}\) in \(\mathfrak{g}\) contains \(\mathcal{C}^k\mathfrak{g} + \mathfrak{h}\) .

